\documentclass[a4paper, 14pt]{article}
% \documentclass[a4paper, 14pt]{extarticle} 

\usepackage[utf8]{inputenc}
\usepackage[T2A]{fontenc}
\usepackage[english,russian]{babel}
% \fontsize{14}{17}

% \usepackage[none]{hyphenat}  % Полное отключение переносов

\usepackage{setspace}
% \onehalfspacing
\usepackage[left=30mm, top=20mm, right=15mm, bottom=20mm, nohead, nofoot]{geometry}
\usepackage{amsmath,amsfonts,amssymb, gensymb, tikz, graphicx}
\usepackage{fancybox,fancyhdr}
\usepackage{multicol}
\usepackage{indentfirst}
\usepackage{titlesec}
\usepackage{listings}
\usepackage{fancyhdr}
\usepackage{float}
\usepackage{tocloft} % для настройки оглавления
\usepackage{etoolbox}
\usepackage{hyperref} % переносит длинные ссылки
\pagestyle{fancy}
\fancyhf{}  % Очистить стандартные колонтитулы
\renewcommand{\headrulewidth}{0pt} % Убираем линию под верхним колонтитулом
\renewcommand{\footrulewidth}{0pt} % Убираем линию над нижним колонтитулом (если нужно)
\cfoot{\thepage}  % Номер справа внизу
% \pagestyle{fancy}
% \fancyhf{}
% \fancyhead[R]{A. Khalitov}
% \fancyfoot[R]{\thepage}
% \fancyhead[L]{Срез знаний за 7 класс}
%\usepackage[ruled,vlined]{algorithm2e}
%\usepackage{algorithm, algorithmic, algpseudocode}
\setcounter{page}{1}
% \headsep=10mm
\usepackage{xcolor}
% \usepackage[style=gost-numeric]{biblatex} % для оформленния ссылок на литературу по ГОСТу
\usepackage{hyperref}
\hypersetup{colorlinks=true, allcolors=black}
\newcommand{\lr}[1]{\left( {#1} \right)}

% Включаем отступ после \subsection
\makeatletter
\patchcmd{\@startsection}{\@afterindentfalse}{\@afterindenttrue}{}{}
\makeatother

%\usepackage{csquotes} % требуется biblatex
%\usepackage[style=numeric,sorting=none]{biblatex}
%\addbibresource{sources.bib}

% Настройка — оставить только название источника в списке литературы
%\AtEveryBibitem{%
%  \clearfield{author}%
%  \clearfield{editor}%
%  \clearfield{publisher}%
%  \clearfield{year}%
%  \clearfield{journal}%
%  \clearfield{note}%
%}

% Переопределим шаблоны, чтобы показывать только заголовок
%\DeclareFieldFormat{title}{#1}
%\DeclareBibliographyDriver{book}{\printfield{title}}
%\DeclareBibliographyDriver{article}{\printfield{title}}
%\DeclareBibliographyDriver{online}{\printfield{title}}

% \usepackage[backend=biber, style=numeric, sorting=none]{biblatex}

% % Настраиваем отображение записей в библиографии
% \DeclareFieldFormat[article,book]{title}{#1} % Название выводится первым
% \renewbibmacro{in:}{} % Убираем "In:" для статей

% % Настраиваем отображение ссылок
% \DeclareFieldFormat{labelnumberwidth}{[#1]} % Формат ссылок: [1], [2], ...

% Отступы
% Абзацный отступ
\setlength{\parindent}{1.25cm}
% Настройка отступов для section
\titleformat{\section}
    {\normalfont\Large\bfseries}{\thesection}{1em}{} % {формат}{метка}{отступ от метки}{перед заголовком}[после заголовка]
\titlespacing*{\section}{1.25cm}{3.5ex}{2.3ex} % {левый отступ}{отступ сверху}{отступ снизу}

% Настройка отступов для subsection
\titleformat{\subsection}
    {\normalfont\large\bfseries}{\thesubsection}{1em}{}
\titlespacing*{\subsection}{1.25cm}{2.5ex}{1.5ex}

% Настройка отступов для subsubsection
\titleformat{\subsubsection}
    {\normalfont\large\bfseries}{\thesubsubsection}{1em}{}
\titlespacing*{\subsubsection}{1.25cm}{2.5ex}{1.5ex}
% \titleformat{\section}{\normalfont\normalsize\bfseries}{\thesection}{1em}{}
% \titleformat{\subsection}{\normalfont\normalsize\bfseries}{\thesubsection}{1em}{}

% Настройка содержания
\renewcommand{\cfttoctitlefont}{\bfseries\fontsize{14}{16}\MakeUppercase}
\renewcommand{\cftaftertoctitle}{\centering} % выравнивание заголовка

% Настройка отступов и стиля пунктов
\renewcommand{\cftsecindent}{0pt} % отступ для разделов
\renewcommand{\cftsubsecindent}{0cm} % отступ для подразделов
\renewcommand{\cftsubsubsecindent}{0cm} % отступ для пунктов

% Настройка точек и шрифта
\renewcommand{\cftsecfont}{\normalfont} % шрифт разделов
\renewcommand{\cftsubsecfont}{\normalfont} % шрифт подразделов
\renewcommand{\cftsubsubsecfont}{\normalfont} % шрифт пунктов

\renewcommand{\cftsecleader}{\cftdotfill{\cftdotsep}} % точки для разделов
\renewcommand{\cftsubsecleader}{\cftdotfill{\cftdotsep}} % точки для подразделов
\renewcommand{\cftsubsubsecleader}{\cftdotfill{\cftdotsep}} % точки для пунктов

% Общий интервал между всеми пунктами
\setlength{\cftbeforetoctitleskip}{0pt} % отступ перед заголовком "Содержание"
\setlength{\cftaftertoctitleskip}{0pt}  % отступ после заголовка "Содержание"
% Расстояние между разделами (section)
\setlength{\cftbeforesecskip}{0pt} % где X - нужное значение в пунктах
% Расстояние между подразделами (subsection)
\setlength{\cftbeforesubsecskip}{0pt}
% Расстояние между подподразделами (subsubsection)
\setlength{\cftbeforesubsubsecskip}{0pt}

% Настройка номеров страниц
\renewcommand{\cftsecpagefont}{\normalfont}   % для разделов (section)
\renewcommand{\cftsubsecpagefont}{\normalfont} % для подразделов (subsection)
\renewcommand{\cftsubsubsecpagefont}{\normalfont} % для пунктов (subsubsection)

% Для листинга кода
\lstset{
  language=python,
  basicstyle=\ttfamily\small,
  keywordstyle=\color{blue}\bfseries,
  stringstyle=\color{red},
  commentstyle=\color{green!50!black}\textit,
  % numbers=left,
  % numberstyle=\tiny\color{gray},
  % stepnumber=1,
  % numbersep=5pt,
  showspaces=false,
  showstringspaces=false,
  % frame=single,
  tabsize=2,
  breaklines=true,
  backgroundcolor=\color{white},
}

\usepackage{pgfplots}
\pgfplotsset{compat=1.18}


% \bibliographystyle{gost-numeric.bbx}
% \usepackage[parentracker=true,
% backend=biber,
% hyperref=false,
% bibencoding=utf8,
% style=numeric-comp,
% language=auto,
% autolang=other,
% citestyle=gost-numeric,
% defernumbers=true,
% bibstyle=gost-numeric,
% sorting=ntvy,
% ]{biblatex}
% \addbibresource{library.bib}

% \documentclass[a4paper,12pt]{article}
% \usepackage[utf8]{inputenc}
% \usepackage[russian]{babel}
% \usepackage{amsmath,amssymb}
% \usepackage{times}
% \usepackage{setspace}
% \usepackage{graphicx}
% \usepackage{float}
% \fontsize{14}{17}
% \usepackage{geometry}
% \usepackage{etoolbox} % Для модификации команд
% \setlength{\parindent}{2cm} % Устанавливаем отступ в преамбуле

% \geometry{
%   left=1.5cm,
%   right=1.5cm,
%   top=2cm,
%   bottom=2cm
% }

% % Включаем отступ после \subsection
% \makeatletter
% \patchcmd{\@startsection}{\@afterindentfalse}{\@afterindenttrue}{}{}
% \makeatother
